\BourbakiExerciseGroup

\begin{enumerate}
\def\labelenumi{\arabic{enumi})}
\item
  假设 A 是一个整环；设 \(f\) 为环 \(\mathbf{A}[X_1, X_2, \ldots, X_n]\) 中次数 \(\leqslant k_i\) （关于 \(X_i\) ，对 \(1 \leqslant i \leqslant n\) ）的一个多项式。对每个 \(i\) （ \(1 \leqslant i \leqslant n\) ）的值，设 \(H_i\) 为 A 的一个含 \(k_i + 1\) 个元素的集合。证明：若 \(f(x_1, x_2, \ldots, x_n) = 0\) 对每个元素 \((x_i) \in \prod_{i=1}^{n} H_i\) 成立，则 \(f = 0\) 。
\item
  假设 A 是一个无限整环；设 \(\Phi\) 为一组多项式 \(\neq0\) （在环 \(A[X_{1},X_{2},\ldots,X_{n}]\) 中）。证明：若 \(\Phi\) 的基数严格小于 A 的基数，则存在 \(A^{n}\) 的一个与 A 等势的子集 H，使得对每个 \(x=(x_{i})\in H\) 及每个 \(f\in\Phi\) ，有 \(f(x)\neq0\) 。
\item
  设 A 是一个整环，并设 B 为 A 的一个无限子集。证明：若多项式 \(f \in A[X]\) 的次数 \textgreater0 ，则 B 在多项式函数 \(x \mapsto f(x)\) 下的像与 B 等势。
\item
  假设 A 的加法群包含一个无限子群 G，其非零元素在 A 中不是零因子。证明映射 \(f \mapsto \tilde{f}\) ——它把 \(A{[}X_{1}, \ldots, X_{p}{]}\) 映入 \(A^{p}\) 到 A 的映射代数——是单射（注意，单未定元的一个 n 次多项式在 G 中不可能有多于 n 个不同的根）。特别地，当 A 中存在一个元素 \(x_{0}\) ——它不是零因子且在 A 的加法群中是无限阶的——时，此结论成立。
\item
  设 K 是一个具有无穷多个元素的域，并设 Q 为 K 上的四元数代数，类型为 \((-1,0,-1)\) （III，第445页）。证明多项式 \(X^{2}+1\) 在 Q 中有无穷多个根。
\item
  设K是一个含有q个元素的有限域。
\end{enumerate}

\begin{enumerate}
\def\labelenumi{\alph{enumi})}
\tightlist
\item
  在环 \(\mathbf{K}[\mathbf{X}_1, \mathbf{X}_2, \ldots, \mathbf{X}_n]\) 中，令 \(\mathfrak{a}\) 为由 \(n\) 个多项式
\end{enumerate}

\(X_{i}^{q}-X_{i}\quad(1\leqslant i\leqslant n)\) 所生成的理想。证明若 \(f\in\mathfrak{a}\) ，则 \(f(x_{1},x_{2},\ldots,x_{n})=0\) 对所有 \((x_{i})\in\mathbf{K}^{n}\) 成立（注意乘法群 \(K^{*}\) 的阶为 q-1）。

\begin{enumerate}
\def\labelenumi{\alph{enumi})}
\setcounter{enumi}{1}
\item
  若 \(f\) 为 \(\mathbf{K}[\mathbf{X}_1, \mathbf{X}_2, \ldots, \mathbf{X}_n]\) 中任意多项式，证明存在恰好一个多项式 \(\vec{f}\) ，使得 \(\deg_i \vec{f} \leqslant q - 1\) （对 \(1 \leqslant i \leqslant n\) ），且 \(f \equiv \vec{f} \pmod{\mathfrak{a}}\) 。于是有 \(\deg \vec{f} \leqslant \deg f\) ；若 \(f\) 使得 \(f(x_1, x_2, \ldots, x_n) = 0\) （对每个 \((x_i) \in \mathbf{K}^n\) ），则 \(f\) 属于理想 \(\mathfrak{a}\) ，因而该理想是同态 \(f \mapsto \vec{f}\) 下 0 的原像（利用习题1）。
\item
  设 \(f_1, f_2, \ldots, f_m\) 为 \(\mathbf{K}[\mathbf{X}_1, \mathbf{X}_2, \ldots, \mathbf{X}_n]\) 中的非零多项式，使得 \(f_i(0, 0, \ldots, 0) = 0 (1 \leqslant i \leqslant m)\) ，且这些 \(f_i\) 的总次数之和 \(< n\) 。证明存在一个元素 \((x_1, x_2, \ldots, x_n) \in \mathbf{K}^n\) ，它异于 \((0, 0, \ldots, 0)\) ，使得
\end{enumerate}

\[
f _ {i} (x _ {1}, x _ {2}, \dots , x _ {n}) = 0 \quad \text { for } \quad 1 \leqslant i \leqslant m.
\]

（观察到如果并非如此，则多项式 \(\prod_{i=1}^{m} (1 - f_i^{q-1})\) 将模 a 与多项式 \(\prod_{i=1}^{n} (1 - X_i^{q-1})\) 同余，并利用 \(b\) ).

\begin{enumerate}
\def\labelenumi{\arabic{enumi})}
\setcounter{enumi}{6}
\item
  设 K 为含 q 个元素的有限域， B 为乘积环 \(K^{1}\) ，其中 I 是一个无限集。给出 B{[}X{]} 中一个非零多项式 f 的例子，使得 \(f(x)=0\) 对所有 \(x\in B\) 成立。
\item
  设 B 为加法群 \(\mathbf{Z} \oplus (\mathbf{Z}/2\mathbf{Z})\) ，并赋予它由令 \((a, b) \cdot (a', b') = (aa', ab' + ba')\) （对 \((a, b) \in \mathbf{B}\) 与 \((a', b') \in \mathbf{B}\) ）所得到的交换环结构。设 \(\varepsilon = (0, 1) \in \mathbf{B}\) ，以及 \(f(\mathbf{X}) = \varepsilon \mathbf{X}(\mathbf{X} - 1) \in \mathbf{B}[\mathbf{X}]\) 。则 \(f(\alpha) = 0\) 对所有 \(\alpha \in B\) 成立。
\item
  设 \(k\) 为整数 \(> 0\) ，使得 \(k \cdot 1 = 0\) （在 \(\mathbf{A}\) 中）。设 \(h\) 为整数 \(> 0\) ， \(f\) 为多项式 \((\mathbf{X} - \alpha)^{k + h} + (\mathbf{X} - \alpha)^k\) 。则 \(\alpha\) 是阶为 \(k\) 的根（对 \(f\) 而言），且是阶 \(\geqslant k + h - 1\) 的根（对 \(Df\) 而言）。
\end{enumerate}

10) 设 \(f\) 为 \(\mathbf{Z}[\mathbf{X}_1, \ldots, \mathbf{X}_n]\) 的一个非零元素。设 \(\mathrm{N}(q)\) 为 \(f\) 在胞腔 \(|x_1| \leqslant q, \ldots, |x_n| \leqslant q\) 中的零点个数。则 \(\mathrm{N}(q) / q^n\) 当 \(q \to +\infty\) 时趋于 0 。

\begin{enumerate}
\def\labelenumi{\arabic{enumi})}
\setcounter{enumi}{10}
\item
  设 P 为所有素数的集合， \(P'\) 为 P 的一个无限子集，且 \(f \in \mathbf{Z}[(\mathbf{X}_i)]\) 。对每个 \(p \in P'\) ，令 \(f_p\) 为 \((\mathbf{Z}/p\mathbf{Z})[(\mathbf{X}_i)]\) 中由对 f 的系数作用 Z 到 Z/pZ 上的典范同态所得到的元素。假设对每个 \(p \in P'\) 及每个 \(x \in (\mathbf{Z}/p\mathbf{Z})^l\) ，我们有 \(f_p(x) = 0\) ；则 f = 0。
\item
  设 \(\mathbf{M}\) 为一个 A-模，且 \(\alpha_0, \ldots, \alpha_n \in \mathbf{A}\) 使得：对 \(i \neq j\) ，关于 \(\alpha_i - \alpha_j\) 的倍乘同态在 \(\mathbf{M}\) 上是双射。对任意 \(m_0, \ldots, m_n \in M\) ，存在且仅存在一个 \(f \in \mathbf{M}[\mathbf{X}]\) ，使得 \(f(\alpha_i) = m_i\) （对 \(0 \leqslant i \leqslant n\) ）且 \(\deg f \leqslant n\) 。
\item
  设 \(\alpha_{i}\) ( \(1 \leqslant i \leqslant n\) ) 为 \(n\) 个互异元素（属于一个交换域 \(\mathbf{K}\) ）， \(\beta_{i}\) ( \(1 \leqslant i \leqslant n\) ) 与 \(\gamma_{i}\) ( \(1 \leqslant i \leqslant n\) ) 为 \(2n\) 个任意元素（属于 \(\mathbf{K}\) ）。证明：存在且仅存在一个多项式 \(f \in \mathbf{K}[\mathbf{X}]\) ，次数 \(\leqslant 2n - 1\) ，使得 \(f(\alpha_{i}) = \beta_{i}\) 且 \(f'(\alpha_{i}) = \gamma_{i}\) （对 \(1 \leqslant i \leqslant n\) ）（埃尔米特插值公式）（首先考虑 \(2n - 1\) 个（在 \(2n\) 个元素 \(\beta_{i}, \gamma_{i}\) 中）为零的情形）。
\end{enumerate}

